\paragraph{Supplement to Finding A52: preserving packets and taking common covers.}
The following observations concern actual pointed covers. They do not
supply a new small integral annihilator or a universal witness.
Fix a genus-two curve $C/k$ over a finite field of characteristic $p$
and $P\in C(k)$. Suppose finitely many constant, connected,
$P$-split cyclic covers $D_i\to C$, of odd degrees prime to $p$,
have good signed $p$-packets in the sense of A49: the corresponding
character values $L_P(\chi,1)$ are units at the selected prime above $p$.
These finitely many packets can be preserved while applying the
variable-quotient supply argument of A49 again.

\begin{proof}
Let $V_i=J_{D_i}[p](\bar k)$ be the finite geometric etale
$p$-torsion vector space. Arithmetic Frobenius $F_i$ acts invertibly
on it and hence has finite order $T_i$. Choose $T$ divisible by all
$T_i$ and all the old deck-Frobenius orders; the latter are one
under the present constant-deck hypothesis. Frobenius commutes
with the deck action. If $j\equiv1\pmod T$, its action over
$k_j=\mathbf F_{q^j}$ on every $V_i$ is exactly the old action:
$F_i^j=F_i$. Thus its invariant subspace in every old character
packet is unchanged. The character-unit criterion of A49 says
precisely that this invariant subspace is zero on a good packet.
A nonzero invariant $p$-primary element would yield an invariant
element of order $p$, so this also controls the corresponding
rational $p$-primary unit test. No rational abelian subvariety
for an individual $p$-adic packet is asserted.

Enlarge the finite excluded-prime set $S$ in A49's
variable-quotient note to contain all old conductor primes, and
let $M$ be its A42 primary-stabilization modulus. The progression
$j\equiv1\pmod{\operatorname{lcm}(M,T)}$ is infinite and has
$\gcd(j,M)=1$. The fixed-kernel counting and Weil estimates of
that note therefore still supply a new good signed packet over
arbitrarily large such extensions. Its conductor can be chosen
coprime to the old ones. Repeating this argument preserves any
finite collection of the actual packets already obtained.
It does not control every faithful packet of one conductor.
\end{proof}

Here is the exact common-cover calculation. Over a common field
where the preceding packets have been preserved, take the connected
component of the fibre product through the selected rational lifts
of $P$. It is an actual connected etale cover $D\to C$ with
constant abelian deck group $H$, a subgroup of the product of the
old cyclic groups. A hypothetical point-section collision gives a
compatible rational lift of $Q$ on this same cover. Put
$\delta=[Q'-P']\in J_D(k)[p^\infty]$. For a quotient
$h_i:D\to D_i$ with kernel $K_i$, actual divisor pushforward and
pullback give
\[
 (h_i)_*\delta=\delta_i,\qquad
 h_i^*(h_i)_*=N_{K_i}:=\sum_{g\in K_i}g.
\tag{CP.1}
\]
Over a finite unramified extension of $\mathbf Z_p$ containing
the roots of unity of order dividing $\exp H$, character evaluation
therefore gives
\[
 \eta(N_{K_i})=
 \begin{cases}
 |K_i|,&\eta|_{K_i}=1,\\
 0,&\eta|_{K_i}\ne1.
 \end{cases}
\tag{CP.2}
\]
Indeed the first sum has all terms one; in the second case
multiplying the sum by a nontrivial value of the character
permutes its terms and forces it to vanish. Since $p\nmid|H|$,
these evaluations describe the full split group algebra and
its action on the actual $p$-primary group.

If an integral operator $B_i$ annihilates $\delta_i$, then
$N_{K_i}\widetilde B_i$ annihilates $\delta$, where
$\widetilde B_i$ is any lift to $\mathbf Z[H]$.
Every integral combination of these lifted relations has zero
character value on a character factoring through none of the
selected quotients. For coprime cyclic conductors this is exactly
the faithful product-conductor limitation of A51.

The genuinely noncyclic test $H=(\mathbf Z/\ell)^2$, with
$\ell\ne p$ odd, gives a more precise boundary. Include all
$\ell+1$ cyclic line quotients. Every nontrivial character $\eta$
has a unique line kernel $K_\eta$, and (CP.2) becomes
\[
 \eta(N_K)=\ell\quad(K=K_\eta),\qquad
 \eta(N_K)=0\quad(K\ne K_\eta).
\tag{CP.3}
\]
All $\eta^u$, $u\in(\mathbf Z/\ell)^\times$, have this same
kernel and form the rational cyclotomic block of that one
quotient. Its signed $p$-orbits may be smaller. Relations lifted
from every other line quotient are zero on the whole block.
Thus accumulating a good packet from each of several different
quotients does not, through these norm identities, fill a missing
signed orbit in any one block. Relabelling a quotient merely
relabells its physical characters; deck translations act within
character summands and do not permute the blocks.

These are identities of actual geometric correspondences, not a
claim that an arbitrary undetected torsion element can be a
compatible point difference. Integral combinations can cancel
coefficients, and no lower bound against every such cancellation
is claimed here. A further geometric relation using the
point-supported nature of $\delta$ could add information, but
none has been obtained in this argument. Preserving finitely many
good packets therefore leaves the universal whole-packet or
small-annihilator supply condition unresolved.
